\chapter{Electrónica Aplicada II: Sensado de Corriente y Finales de Carrera Electrónicos}
\label{sec:electronica_aplicada_2}

En el marco de la asignatura Electrónica Aplicada II, el proyecto aborda el diseño de una etapa analógica de alta
precisión para la medición continua de corriente en los motores de accionamiento y la detección electrónica de topes
mecánicos (\textit{stall detection}).

\section{Justificación Técnica y Objetivos del Módulo Analógico}
\label{sec:ea2_objetivos}

Los sistemas CNC convencionales emplean interruptores electromecánicos o sensores ópticos como finales de carrera en
los extremos de cada eje. En el presente proyecto, con el propósito de maximizar el aprovechamiento de los conceptos
de amplificadores operacionales y procesamiento analógico de señales, se implementa un sistema de final de carrera
electrónico basado en la detección del incremento severo de corriente cuando un eje alcanza su límite físico o sufre
un atascamiento mecánico.

El módulo analógico debe cumplir con las siguientes funciones primarias:
\begin{itemize}
  \item Muestreo de la caída de tensión en una resistencia de valor muy bajo ($R_{shunt}$) ubicada en la rama de retorno
    de cada canal del controlador de potencia.
  \item Amplificación diferencial con rechazo en modo común para eliminar el ruido provocado por las conmutaciones de
    potencia.
  \item Filtrado pasa-bajos para atenuar componentes armónicas de alta frecuencia generadas por la modulación por ancho
    de pulso (PWM).
  \item Acondicionamiento de señal con doble salida: una señal analógica continua proporcional a la corriente del
    motor dirigida al conversor analógico-digital (ADC) y una salida digital producida por un comparador con
    histéresis para disparo inmediato de interrupción.
\end{itemize}

\section{Diseño del Circuito de Sensado con Amplificadores Operacionales}
\label{sec:ea2_circuito}

En la figura~\ref{crkt.ej2:sensado_opamp} se ilustra el esquema eléctrico propuesto para la etapa de acondicionamiento.
La tensión en extremos de la resistencia derivadora $R_{shunt}$ es tomada por un amplificador operacional en
configuración diferencial.

\begin{figure}[!ht]
  \centering
  \begin{tikzpicture}[
  >=stealth, thick,
  font=\small,
  node distance=1.2cm and 1.8cm,
  block/.style={draw=blue!70!black, fill=blue!10, rectangle, rounded corners=4pt, minimum height=1.1cm, minimum width=2.5cm, align=center},
  inputnode/.style={draw=gray!70!black, fill=gray!12, rectangle, rounded corners=4pt, minimum height=1.1cm, minimum width=2.4cm, align=center},
  outnode/.style={draw=green!60!black, fill=green!12, rectangle, rounded corners=4pt, minimum height=1.1cm, minimum width=2.6cm, align=center}
]

  % Fila 1: Cadena Analógica Principal
  \node[inputnode] (trans) {Sensado\\{\footnotesize ($I_{motor}$ / Shunt)}};
  \node[block, right=of trans] (amp) {Ganancia y\\Acondic.\\{\footnotesize (Diferencial)}};
  \node[block, right=of amp] (filtro) {Filtrado\\Pasa-Bajos\\{\footnotesize (Rechazo PWM)}};
  \node[outnode, right=of filtro] (adc) {Salida Analógica\\{\footnotesize ($V_{ana} \rightarrow$ ADC MCU)}};

  % Fila 2: Etapa de Comparación y Disparo
  \node[block, below=1.3cm of filtro] (comp) {Comparador con\\Histéresis\\{\footnotesize (Schmitt Trigger)}};
  \node[outnode, right=of comp] (exti) {Disparo Digital\\{\footnotesize (EXTI / Tope)}};

  % Conexiones Horizontales con etiquetas arriba
  \draw[->] (trans) -- node[above, font=\footnotesize] {$V_{sens}$} (amp);
  \draw[->] (amp) -- node[above, font=\footnotesize] {$V_{amp}$} (filtro);
  \draw[->] (filtro) -- node[above, font=\footnotesize] {$V_{ana}$} (adc);

  % Conexión vertical directa hacia la etapa comparadora
  \draw[->] (filtro) -- node[right, font=\footnotesize] {$V_{ana}$} (comp);
  \draw[->] (comp) -- (exti);

  % Entrada de umbral de referencia
  \draw[<-] (comp.south) -- ++(0,-0.7) node[below, font=\footnotesize] {Umbral Ref. ($V_{ref}$)};

\end{tikzpicture}

  \caption{Esquema circuital de sensado de corriente con amplificador diferencial, filtro RC y comparador Schmitt.}
  \label{crkt.ej2:sensado_opamp}
\end{figure}

La relación de transferencia teórica para la etapa amplificadora diferencial está dada por la ecuación~\ref{eq.ej2:vo_diff}:

\begin{equation}
  V_{diff} = \left( \frac{R_f}{R_1} \right) \cdot I_{motor} \cdot R_{shunt}
  \label{eq.ej2:vo_diff}
\end{equation}

Donde $R_f / R_1$ establece la ganancia de tensión ($A_v$) requerida para escalar la militensión en el shunt hacia el
rango dinámico del ADC del microcontrolador ($0\V$ a $3.3\V$).

\section{Filtrado y Detector de Umbral con Histéresis (Schmitt Trigger)}
\label{sec:ea2_schmitt}

La corriente en los devanados del motor presenta rizo debido a la conmutación de los transistores del driver. Para
evitar falsos disparos en el detector de límite, la señal amplificada atraviesa un filtro pasa-bajos pasivo de primer
orden ($R_{flt}, C_{flt}$) cuya frecuencia de corte se define en la ecuación~\ref{eq.ej2:fc_filter}:

\begin{equation}
  f_c = \frac{1}{2 \pi \cdot R_{flt} \cdot C_{flt}}
  \label{eq.ej2:fc_filter}
\end{equation}

La señal filtrada $V_{ana}$ se aplica a la entrada inversora de un comparador configurado con histéresis no inversora o
inversora (disparador Schmitt). El umbral de referencia $V_{ref}$ se ajusta mediante un divisor resistivo con
potenciómetro, fijando el nivel de corriente equivalente al atascamiento físico del cabezal.

Cuando $V_{ana}$ supera a $V_{ref}$, la salida del comparador conmuta de estado lógico alto a bajo, generando una
transición limpia requerida por los pines de interrupción del controlador. La histéresis añadida mediante $R_{h1}$ y
$R_{h2}$ previene la oscilación indeseada en la zona de transición.

\section{Flujo de Operación del Sensado}
\label{sec:ea2_flujo}

La figura~\ref{fig.ej2:flujo_sensado} representa el algoritmo analógico-digital de supervisión de la corriente de los
motores.

\begin{figure}[!ht]
  \centering
  \begin{tikzpicture}[node distance=1.6cm, auto, >=stealth, thick,
  startstop/.style={draw, fill=red!15, rectangle, rounded corners, minimum height=0.8cm, text centered, align=center},
  process/.style={draw, fill=blue!10, rectangle, minimum height=1.0cm, text centered, align=center},
  decision/.style={draw, fill=yellow!15, diamond, aspect=2, inner sep=2pt, text centered, align=center}]

  \node [startstop] (start) {Inicio Muestreo de Corriente};
  \node [process, below=of start] (read) {Lectura de Caída en Shunt ($V_{shunt}$)};
  \node [process, below=of read] (amp) {Amplificación Diferencial y Filtrado RC};
  \node [decision, below=1.2cm of amp] (check) {$V_{ana} > V_{ref}$? \\ (¿Tope de Eje?)};
  \node [process, right=2.0cm of check] (adc) {Enviar $V_{ana}$ a ADC \\ (Monitoreo Carga)};
  \node [process, below=1.4cm of check] (interrupt) {Disparo Comparador Schmitt \\ Interrupción EXTI};
  \node [startstop, below=of interrupt] (stop) {Detener Motores PaP \\ (Parada Emergencia)};

  \draw [->] (start) -- (read);
  \draw [->] (read) -- (amp);
  \draw [->] (amp) -- (check);
  \draw [->] (check) -- node[above, font=\small] {No} (adc);
  \draw [->] (check) -- node[right, font=\small] {Sí} (interrupt);
  \draw [->] (interrupt) -- (stop);
  \draw [->] (adc.north) |- (read.east);
\end{tikzpicture}

  \caption{Diagrama de flujo de procesamiento de la señal de corriente y disparo de final de carrera.}
  \label{fig.ej2:flujo_sensado}
\end{figure}

Este esquema permite reemplazar la necesidad de finales de carrera mecánicos tradicionales por una solución
electrónica integrada en las placas de control.
